\chapter*{Résumé}
\addcontentsline{toc}{chapter}{Résumé}

Le présent rapport rend compte du stage de fin d'études effectué au sein de l'ENSA Béni Mellal, rattachée à l'Université Sultan Moulay Slimane. Ce stage s'inscrit dans le cadre de l'obtention du Bachelor Universitaire de Technologie en Génie Logiciel de l'EST de Fès.

Le projet consiste en la conception et le développement d'une plateforme web intelligente dédiée à la présélection des candidats aux filières d'ingénieur de l'ENSA BM. Cette plateforme automatise le traitement des dossiers de candidature, depuis le dépôt en ligne jusqu'au classement final, en passant par la vérification OCR des documents, le scoring académique pondéré par filière et la gestion des épreuves écrites.

La solution repose sur une architecture Full Stack moderne : Django REST Framework pour le backend et React 18 avec Tailwind CSS pour le frontend. Elle intègre un moteur de règles de rejet automatique (7 types), un algorithme de scoring avec fuzzy matching des matières, un système de notifications en temps réel via WebSocket, et un module d'import des notes d'épreuve écrite au format Excel.

\textbf{Mots-clés :} Présélection, ENSA, Django, React, Scoring automatique, OCR, Fuzzy Matching, WebSocket, LMD, Full Stack.
